\chapter{PROTOTIPAR}
El presente capítulo desarrolla la cuarta etapa del Design Thinking: Prototipar. Su propósito es construir una representación tangible de las ideas seleccionadas en el Capítulo 6, que permita visualizar las decisiones de diseño, iniciar conversaciones productivas con los usuarios y detectar errores de forma temprana (Vargas Márquez et al., 2021). Se describe la arquitectura del sistema, las reglas de interacción implementadas, el algoritmo de emparejamiento, el sistema de diseño bajo principios de IHC Afectiva, y las iteraciones realizadas durante el desarrollo.

[Imagen: mockup general de la aplicación ECCO, mostrando las pantallas principales en secuencia]

\section{ARQUITECTURA}
El prototipo se desarrolló bajo una arquitectura cliente-servidor, compuesta por los siguientes elementos:

\begin{itemize}
    \item \textbf{Cliente:} aplicación desarrollada en React Native con Expo (TypeScript), organizada por funcionalidades: autenticación, \textit{onboarding}, inicio, conversación, jardín, perfil y ajustes. El uso de Expo permite probar la aplicación directamente en el teléfono mediante Expo Go, sin necesidad de compilación previa, lo cual facilita la ejecución de las pruebas piloto descritas en el Capítulo 8.
    \item \textbf{Servidor:} API REST desarrollada en FastAPI (Python), organizada en los módulos de autenticación, usuarios, ánimo, emparejamiento, chat, jardín y sugerencias.
    \item \textbf{Base de datos:} PostgreSQL, con migraciones versionadas mediante Alembic. Se optó por un modelo relacional en lugar de uno documental, dado que los datos del sistema (usuarios, conversaciones, mensajes y decisiones) presentan relaciones estructuradas entre sí, y porque PostgreSQL permite bloquear el proceso de emparejamiento para evitar que dos búsquedas simultáneas generen parejas duplicadas.
    \item \textbf{Servicio externo:} Google Knowledge Graph Search API, consultada únicamente desde el servidor para proteger la clave de acceso, utilizada para el autocompletado de gustos específicos (Google, s.f.; sección 7.4).
\end{itemize}

[Imagen: diagrama de arquitectura cliente-servidor de ECCO]

La aplicación consulta periódicamente al servidor para recibir mensajes nuevos y para verificar si ya se formó una pareja (cada 4 segundos dentro de una conversación activa, y cada 3 a 5 segundos durante la búsqueda), y únicamente mientras la pantalla correspondiente está visible, con el fin de equilibrar la capacidad de respuesta percibida con el consumo de recursos del dispositivo.

En cuanto a la calidad del software, el servidor cuenta con 45 pruebas automatizadas que verifican los módulos de autenticación, perfil, emparejamiento, chat, reglas semanales, jardín, reportes y sugerencias. Adicionalmente, el flujo completo de la aplicación fue probado con dos usuarios simulados en entorno de navegador.

\section{REGLAS DE LA INTERACCIÓN}
La Tabla \ref{tab:reglas_interaccion} resume los parámetros que rigen la interacción dentro del prototipo. Estos valores constituyen hipótesis de diseño derivadas de la literatura y de la evidencia primaria recolectada en los Capítulos 4 y 5, y se someterán a validación durante la prueba piloto (Capítulo 8).

\begin{longtable}{
  >{\raggedright\arraybackslash}p{0.55\textwidth}
  >{\raggedright\arraybackslash}p{0.35\textwidth}
}
\caption{Reglas de la interacción}\label{tab:reglas_interaccion}\\
\toprule
\textbf{Regla} & \textbf{Valor} \\
\midrule
\endfirsthead
\multicolumn{2}{l}{\textit{Tabla \thetable{} (continuación)}} \\[4pt]
\toprule
\textbf{Regla} & \textbf{Valor} \\
\midrule
\endhead
\midrule
\multicolumn{2}{r}{\textit{Continúa en la siguiente página}} \\
\endfoot
\bottomrule
\endlastfoot

Duración de una conversación & 7 días \\
Días de gracia para decidir & 3 (del día 7 al 10) \\
Mensajes por persona y por día & 5 \\
Largo máximo de un mensaje & 500 caracteres \\
Etapas de la planta & 6: semilla, brote, plantita, capullo, primera flor y flor completa. Avanza una etapa por cada día en que ambos usuarios escriben \\
Ánimo visible para la pareja & Solo el de las últimas 24 horas; puede desactivarse \\
Duración máxima de una búsqueda & Hasta 24 horas; el usuario puede salir y la búsqueda continúa en segundo plano \\
Temas por búsqueda & Hasta 3, seleccionados entre los propios intereses del usuario \\
Intereses mínimos requeridos & 3 \\

\end{longtable}

El límite de cinco mensajes diarios busca que cada mensaje sea redactado y dirigido personalmente a la otra persona, el tipo de comunicación que Burke y Kraut (2016) asociaron con mejoras en el bienestar, y evita a la vez incorporar elementos orientados a prolongar el tiempo de uso (Montag et al., 2019), en coherencia con la asociación entre el uso problemático de las redes sociales y una mayor soledad (Feng et al., 2026). Este valor constituye una hipótesis de diseño que se validará en la prueba piloto (Capítulo 8).

\section{ALGORITMO DE EMPAREJAMIENTO}
Cuando un usuario inicia una búsqueda, el sistema lo compara con todas las personas que se encuentran buscando pareja en ese mismo momento, y selecciona a aquella con el puntaje de afinidad más alto, calculado mediante la siguiente fórmula:

\begin{quote}
\textbf{Puntaje = 3 × temas en común + 1 × intereses en común + 2 × subtemas en común + 4 × ejemplos en común}
\end{quote}

Esta fórmula otorga mayor peso a los niveles más específicos de afinidad: compartir un ejemplo concreto (por ejemplo, la misma película) pesa más que compartir únicamente la categoría general (por ejemplo, ``Películas''). Este diseño responde directamente a los hallazgos de los Capítulos 4 y 5, según los cuales los usuarios valoran la posibilidad de conectar a partir de gustos específicos y particulares, y es coherente con el meta-análisis de Montoya et al. (2008), según el cual la similitud real entre las personas se asocia con la atracción interpersonal, especialmente en interacciones breves o iniciales.

La comparación de textos ignora mayúsculas, espacios y tildes (por ejemplo, ``Iron Man 2'' equivale a ``iron man 2''). El sistema garantiza que nunca se repita una pareja ya formada previamente, y un bloqueo a nivel de base de datos evita que dos búsquedas simultáneas generen parejas duplicadas.

[Imagen: diagrama de flujo del algoritmo de emparejamiento]

\section{JERARQUÍA DE INTERESES Y AUTOCOMPLETADO}
El catálogo de intereses del prototipo está organizado en 11 categorías y 77 intereses: Entretenimiento, Música, Cultura, Cocina y gastronomía, Bebidas, Deportes, Actividades al aire libre, Estilo de vida, Bienestar, Creatividad, y Social y comunidad. Cada interés cuenta con sus propios subtemas (por ejemplo, Películas tiene géneros, Rock tiene estilos musicales, y Fútbol tiene ligas), y dentro de cada subtema el usuario puede agregar ejemplos concretos de forma opcional (Películas → Acción → \textit{Iron Man 2}) o de manera general (Películas → \textit{Coco}).

Para facilitar el ingreso de ejemplos concretos, la aplicación ofrece autocompletado mediante la Google Knowledge Graph Search API, que sugiere nombres reales filtrados según el tipo de interés (en Películas sugiere únicamente películas; en Rock, bandas y artistas; en Fútbol, equipos y jugadores). El usuario conserva siempre la posibilidad de escribir texto libre si no encuentra el ejemplo que busca.

Dado que la API externa responde en aproximadamente un segundo, lo cual excede el Umbral de Doherty de 400 ms que optimiza la productividad percibida (Yablonski, 2020), la interfaz muestra de inmediato las sugerencias ya recibidas que coinciden parcialmente con lo que el usuario va escribiendo, mientras el servidor almacena en caché las respuestas por 24 horas y reutiliza la conexión con el servicio externo, reduciendo así la latencia percibida en búsquedas posteriores.

[Imagen: pantalla de selección de intereses con autocompletado en acción]

\section{SEGURIDAD Y PRIVACIDAD}
Dado que la aplicación empareja a jóvenes desconocidos entre sí para conversar, se implementaron los siguientes mecanismos de seguridad y protección de la privacidad, en respuesta directa a la necesidad de confianza identificada en la Síntesis de la Fase Empatizar (4.4):

\begin{itemize}
    \item Las contraseñas se almacenan mediante hash Argon2, y se exige un mínimo de 8 caracteres, una mayúscula, un número y un símbolo.
    \item La sesión del usuario se gestiona mediante un token JWT válido por 7 días, almacenado en el almacenamiento seguro del dispositivo.
    \item Se establece un límite de intentos de inicio de sesión y de recuperación de contraseña; la recuperación se realiza mediante un enlace enviado por correo electrónico, con una vigencia de 30 minutos.
    \item La aplicación no solicita fotografías del usuario en ningún momento. El estado de ánimo compartido es opcional, y la decisión final de continuidad de una conversación es privada: un ``no'' de la otra persona no se revela mientras el usuario espera el resultado, evitando así que se infiera el voto ajeno a partir del tiempo de respuesta.
    \item El sistema de reporte permite indicar un motivo (me hizo sentir incómodo/a, contenido ofensivo, acoso o presión, \textit{spam}, otro) junto con un comentario opcional. La conversación se cierra de inmediato y la pareja reportada no vuelve a ser emparejada en el futuro.
    \item Los acuerdos de convivencia, presentados en la pantalla ``Un lugar seguro'' del \textit{onboarding}, aclaran explícitamente que ECCO no reemplaza la ayuda profesional, en línea con los límites definidos en la sección 1.5 del presente documento.
\end{itemize}

[Imagen: pantalla ``Un lugar seguro'' con los acuerdos de convivencia]

\section{SISTEMA DE DISEÑO: IHC AFECTIVA Y UX/UI EMOCIONAL}
El sistema de diseño del prototipo se construyó deliberadamente para traducir en elementos visuales e interactivos concretos los principios de IHC Afectiva revisados en el Capítulo 2, evitando los códigos visuales que la literatura del Capítulo 3 asocia a la ansiedad social y la comparación: colores de alerta, contadores públicos y notificaciones de enganche.

\subsection{Paleta de colores y accesibilidad}
La paleta es cálida y natural, compuesta por un color durazno como color principal, verde para representar el crecimiento y las acciones positivas, y celeste para transmitir calma y acceso a la ayuda, complementada con cinco colores adicionales utilizados en la representación de los estados de ánimo.

[Imagen: paleta de colores de ECCO, con los tonos principales y secundarios]

Dado que la calidez cromática por sí sola no garantiza la legibilidad, se midió el contraste de todos los elementos de texto e interactivos conforme a las Pautas de Accesibilidad para el Contenido Web (WCAG) 2.2 (W3C, 2023), y se corrigieron los valores que no alcanzaban el mínimo recomendado, sin abandonar la paleta original. La Tabla \ref{tab:contraste} resume los ajustes realizados.

\begin{longtable}{
  >{\raggedright\arraybackslash}p{0.38\textwidth}
  >{\centering\arraybackslash}p{0.26\textwidth}
  >{\centering\arraybackslash}p{0.26\textwidth}
}
\caption{Corrección de contraste según WCAG 2.2}\label{tab:contraste}\\
\toprule
\textbf{Elemento} & \textbf{Antes} & \textbf{Después} \\
\midrule
\endfirsthead
\multicolumn{3}{l}{\textit{Tabla \thetable{} (continuación)}} \\[4pt]
\toprule
\textbf{Elemento} & \textbf{Antes} & \textbf{Después} \\
\midrule
\endhead
\midrule
\multicolumn{3}{r}{\textit{Continúa en la siguiente página}} \\
\endfoot
\bottomrule
\endlastfoot

Texto blanco sobre botón durazno & 1,95:1 & 5,27:1 (café oscuro sobre durazno) \\
Botones verdes y celestes & 1,86:1 y 1,62:1 & 5,30:1 y 5,71:1 \\
Enlaces & 2,84:1 & 4,72:1 -- 7,17:1 \\
Logo y saludo sobre el fondo verde de Inicio & 2,98:1 & 6,1:1 \\

\end{longtable}

Adicionalmente, todos los objetivos táctiles de la interfaz miden al menos entre 40 y 48 px, conforme a la Ley de Fitts, que establece que el tiempo requerido para alcanzar un objetivo es función de su tamaño y distancia (Yablonski, 2020).

\subsection{Tipografía}
Se utiliza Poppins para los títulos y Quicksand para el cuerpo de texto, con un tamaño mínimo de 12 px en toda la interfaz. Esta combinación tipográfica, de formas redondeadas, refuerza la calidez emocional de la aplicación en coherencia con la faceta de lo deseable propuesta por Morville (2005, como se citó en Méndez et al., 2014), evitando las fuentes geométricas y frías asociadas a interfaces de carácter corporativo.

[Imagen: muestra tipográfica de Poppins y Quicksand aplicadas en una pantalla de la app]

\subsection{Ilustraciones propias y mascota}
Con el fin de construir una identidad visual propia y coherente, se diseñaron ilustraciones originales en formato vectorial (SVG): una semilla de bienvenida, cinco caras que representan los estados de ánimo (fatal, mal, normal, bien, feliz), y un fondo de colinas y arbustos para la pantalla de Inicio, que se adapta a cualquier ancho de pantalla sin deformarse.

[Imagen: ilustración de la semilla de bienvenida]
[Imagen: las cinco caras de ánimo]
[Imagen: fondo de colinas de la pantalla de Inicio]

\subsection{Micro-interacciones}
Las micro-interacciones del prototipo se diseñaron siguiendo la definición de Saffer (2013), quien sostiene que su poder reside en la brevedad: no requieren un compromiso prolongado del usuario para generar un efecto significativo. Entre las implementadas en ECCO se encuentran:

\begin{itemize}
    \item La planta crece con una breve animación acompañada del aviso ``¡Su planta creció!'', ofreciendo retroalimentación inmediata ante un comportamiento recíproco entre ambos usuarios, en línea con los hallazgos de Godard y Holtzman (2024) sobre el valor del uso activo y recíproco.
    \item Una animación de búsqueda con mensajes que cambian progresivamente, reduciendo la percepción de espera mientras el sistema localiza una posible pareja (Umbral de Doherty; Yablonski, 2020).
    \item El avatar realiza un pequeño ``salto'' al ser personalizado, y la cara de ánimo elegida se agranda al momento de registrarse, ofreciendo retroalimentación informativa inmediata ante cada acción del usuario (Shneiderman et al., 2016).
    \item Un ejercicio de respiración guiada, representado mediante un círculo que se expande y se contrae, se activa cuando el usuario registra un estado de ánimo bajo, en línea con el respaldo empírico a los ejercicios guiados en aplicaciones de bienestar (Eisenstadt et al., 2021).
    \item La pantalla de cierre de cada semana fue diseñada con especial cuidado visual y de contenido, dado que, conforme a la Regla del Pico-Final, las personas juzgan una experiencia principalmente por su momento más intenso y por su cierre (Yablonski, 2020).
\end{itemize}

[Imagen: secuencia de la animación de crecimiento de la planta]
[Imagen: pantalla de respiración guiada]
[Imagen: pantalla de cierre de semana]

\subsection{Lenguaje del sistema}
Todos los mensajes generados por el sistema (avisos, confirmaciones, estados de espera) se redactaron en segunda persona, con un tono cálido y libre de culpa. Por ejemplo, al agotarse los mensajes diarios el sistema indica ``Ya compartiste tus mensajes de hoy. Mañana podrás escribir de nuevo'', en lugar de un aviso de error genérico; y la opción de abandonar una conversación se etiqueta como ``Cerrar aquí'' en lugar de ``Mejor no''. Esta decisión responde al enfoque afectivo de la IHC, que no solo busca que la interacción resulte placentera, sino también influir en las emociones y actitudes del usuario respecto al uso del sistema (Gupta, 2012, como se citó en Méndez et al., 2014), y a la recomendación de ofrecer instrucciones sencillas y constructivas cuando el usuario no puede completar una acción (Shneiderman et al., 2016).

\subsection{Diseño adaptable}
La interfaz se adapta a teléfonos pequeños, grandes y tablets. En la pantalla de Inicio, el contenido ocupa todo el ancho disponible y conserva el mismo orden que en la versión para teléfono, evitando el uso de columnas múltiples que fragmenten la atención del usuario, en coherencia con la Ley de Miller sobre la capacidad limitada de la memoria de trabajo (Yablonski, 2020).

\section{ITERACIONES DURANTE EL DESARROLLO}
Durante la construcción del prototipo se identificaron problemas de usabilidad que motivaron ajustes de diseño, lo cual ilustra el carácter iterativo del Design Thinking, que permite retornar a fases anteriores cuando los hallazgos así lo exigen (Vargas Márquez et al., 2021):

\begin{enumerate}
    \item \textbf{Temas de búsqueda:} inicialmente se mostraban seis temas fijos, independientemente de los intereses elegidos por el usuario. Se corrigió para mostrar únicamente los intereses propios de cada usuario.
    \item \textbf{Búsqueda en segundo plano:} no quedaba claro cómo retomar una búsqueda tras salir de la pantalla. La tarjeta de la pantalla de Inicio se actualiza ahora automáticamente y notifica ``¡Encontramos a alguien!'' cuando corresponde.
    \item \textbf{Animación de bienvenida:} se probó inicialmente una animación desarrollada en Rive, la cual no es compatible con Expo Go, herramienta utilizada en las pruebas piloto. Se reemplazó por una ilustración SVG propia.
    \item \textbf{Ancho en pantallas grandes:} se evaluaron un ancho máximo fijo y un diseño de dos columnas; finalmente se optó por una sola columna que ocupa todo el ancho disponible, por su mayor coherencia con la Ley de Miller.
    \item \textbf{Gustos específicos:} se incorporó la jerarquía de subtemas y ejemplos con autocompletado (sección 7.4), con el fin de que el algoritmo de emparejamiento (sección 7.3) pudiera operar con mayor precisión.
    \item \textbf{Retroalimentación del autocompletado:} cuando el servicio externo no encuentra sugerencias, el sistema informa ``Sin sugerencias. Toca + para agregarlo tal como lo escribiste'', atendiendo el principio de retroalimentación, que consiste en comunicar de inmediato al usuario el resultado de sus acciones (Norman, 2013). También se redujo el tiempo de espera percibido mediante el almacenamiento en caché descrito en la sección 7.4.
\end{enumerate}

[Imagen: comparación antes/después de la pantalla de selección de temas]

\section{HACIA LA FASE TESTEAR}
Las decisiones de diseño descritas en este capítulo constituyen un conjunto de hipótesis que deberán validarse con usuarios reales durante la prueba piloto, desarrollada en el Capítulo 8. En particular, interesa conocer si la duración de siete días y el límite de cinco mensajes diarios resultan adecuados, si las ideas para romper el hielo cumplen su función, y si la representación del vínculo mediante una planta motiva la interacción sin generar presión, en coherencia con el principio de poner los prototipos en manos de los usuarios para recopilar retroalimentación cualitativa y cuantitativa (Vargas Márquez et al., 2021).